\documentclass[11pt]{article}
\usepackage[margin=1in]{geometry}
\usepackage{amsmath,amssymb}
\usepackage{booktabs}
\usepackage{hyperref}
\hypersetup{colorlinks=true, linkcolor=blue, urlcolor=blue}

\title{Minimum-Time Low-Thrust Transfer from a Distant Retrograde Orbit\\
       to a Lunar Tulip Orbit\\
       \large Formulation, discretization, and what the first solves taught us}
\author{M.\ Casey}
\date{\today}

\newcommand{\vecr}{\mathbf{r}}
\newcommand{\vecv}{\mathbf{v}}
\newcommand{\bal}{\boldsymbol{\alpha}}
\newcommand{\code}[1]{\texttt{#1}}

\begin{document}
\maketitle

\begin{abstract}
We pose and solve the minimum-time low-thrust transfer between two periodic
orbits of the Earth--Moon circular restricted three-body problem: a distant
retrograde orbit about the Moon, and a seven-petal tulip orbit whose apoapses
hang over the lunar southern hemisphere. The problem is solved by direct
collocation and compared against an existing indirect (Pontryagin shooting)
solution of the same transfer---the first problem in this programme where an
independent second opinion exists. The headline result is a certified
agreement: with fourth-order Hermite--Simpson at $N = 1600$, the direct solve
reproduces the indirect $t_f$ to five significant figures ($4.0152501$ against
$4.0152425$), with a worst-interval position error of $3.3$~m and an end-to-end
propagated position error of $42$~m. That is the direct-against-indirect
cross-validation this programme was built to perform. Reaching it required
discarding the second-order scheme entirely: trapezoidal solutions that
converged to defects of $10^{-14}$, and that appeared to \emph{beat} the
reference minimum time, are not trajectories of the continuous problem at all.
\textbf{A machine-tight defect is a statement about the discretization, not
about the trajectory}---here the two differ by orders of magnitude. We also
report a minimum-altitude path constraint, why adding one costs a block of
algebraic inequalities in the direct formulation against a multipoint
boundary-value problem in the indirect one, and the pass/fail gate that
separates a converged NLP from an answer. The comparison the campaign was built
for is then carried out in full: the direct NLP multipliers, mapped to costates,
reproduce the indirect costate field to $6\times10^{-4}$ degrees in direction
and to a scale factor of $0.99999$, and the state histories agree to $0.4$~km
once a $2.9$-second phasing offset is removed---the two methods find the same
extremal, not merely the same final time. One \emph{negative} result recurs
through the document and is stated plainly rather than buried: we twice
concluded that the unconstrained problem is ill-posed, and twice the evidence
turned out to be discretization error rather than physics. That conjecture
remains open.
\end{abstract}

\tableofcontents

\section{The CR3BP frame and its coordinates}
\label{sec:frame}

Every number in this document lives in one particular reference frame, and the
frame is not inertial. It is worth being precise about it before any orbit is
described, because two of the terms in the equations of motion exist only
because of the frame's rotation.

\paragraph{The construction.} The circular restricted three-body problem
idealizes the Earth and the Moon---the two \emph{primaries}---as moving on
circular orbits about their common barycentre with constant angular rate
$n = \sqrt{G(M_E+M_M)/L^3}$, and the spacecraft as a massless third body that
feels both but perturbs neither. ``Restricted'' is exactly that massless
assumption; ``circular'' is the assumption that the Moon's real eccentricity of
$0.055$ can be ignored.

Now spin the coordinate frame at that same rate $n$. In the rotating frame the
two primaries stand \emph{still}. That is the whole point of the construction:
a problem with two moving attractors becomes a problem with two fixed ones, at
the price of two fictitious accelerations.

\paragraph{Axes and origin.} The origin is the Earth--Moon barycentre. Then

\begin{itemize}
  \item $\hat{x}$ points from the Earth toward the Moon, along the line joining
        the primaries;
  \item $\hat{z}$ is along the orbital angular momentum of the primaries, i.e.
        the normal to their orbit plane---so positive $z$ is lunar north;
  \item $\hat{y} = \hat{z} \times \hat{x}$ completes the right-handed set, and
        points along the Moon's direction of travel.
\end{itemize}

With the mass ratio
\[
  \mu = \frac{M_M}{M_E + M_M} = 0.012150585609624,
\]
the primaries sit permanently at
\[
  \vecr_{\!E} = (-\mu,\,0,\,0), \qquad \vecr_{\!M} = (1-\mu,\,0,\,0),
\]
one nondimensional unit apart. Throughout, $d = \lVert\vecr - \vecr_{\!E}\rVert$
is the distance to the Earth and $\rho = \lVert\vecr - \vecr_{\!M}\rVert$ the
distance to the Moon; the ``lunar altitude'' quoted everywhere below is
$\rho L - R_M$ with $R_M = 1737.4$~km.

\paragraph{What the frame costs.} A frame rotating at rate $n\hat z$ introduces
centrifugal and Coriolis accelerations. Nondimensionalized so that $n = 1$
(Section~\ref{sec:units}), the centrifugal term is $(x, y, 0)$---it appears
folded into $\mathbf{g}$ in \eqref{eq:vdot}, which is why the first two rows of
$\mathbf{g}$ begin with a bare $x$ and $y$---and the Coriolis term is
$\mathbf{h}(\vecv) = (2\dot y,\, -2\dot x,\, 0)$. Neither is a force; both are
bookkeeping for having chosen a rotating frame. They are also the reason the
CR3BP admits the five Lagrange equilibria and the whole zoo of periodic families
that the endpoints of this transfer are drawn from.

\paragraph{A caution about what the pictures show.} A trajectory that looks
like a tidy closed loop in this frame is, in the inertial frame, a rosette; and
a trajectory that appears to loop \emph{backward} around the Moon (the
``retrograde'' in DRO) may be perfectly prograde inertially. All plots and
movies in this document are rotating-frame plots unless stated otherwise.

\section{From dimensional to nondimensional units}
\label{sec:units}

The CR3BP is customarily nondimensionalized so that three quantities equal one:
the distance between the primaries, their total mass, and their angular rate.
This is not cosmetic. It puts positions, velocities and times all within an
order of magnitude of unity, which is what lets an interior-point NLP solver
work at a single tolerance across the whole state vector instead of fighting a
$10^{5}$~km position against a $10^{-3}$~km/s velocity.

\paragraph{The three scale factors.} Everything follows from a length, a time
and a mass:
\[
  L = 389703.264829278~\text{km}, \qquad
  T = 382981.289129055~\text{s}, \qquad
  M = M_E + M_M ,
\]
with $T = \sqrt{L^3 / GM}$ chosen precisely so that the primaries' angular rate
is $n = 1/T$, i.e.\ unity in the new units. These are the values \code{pumpkyn}
uses, and we adopt them verbatim so that our direct solutions and its indirect
ones are solutions of the \emph{same} problem---a shared constant set is a
precondition for the cross-validation this campaign exists to perform.

Two derived scales do most of the work:
\[
  V = \frac{L}{T} = 1.0175517~\text{km/s},
  \qquad
  A = \frac{L}{T^2} = 2.656923~\text{mm/s}^2 .
\]

\paragraph{The conversions.} To go from dimensional (left) to nondimensional
(right):

\begin{center}
\begin{tabular}{llll}
\toprule
Quantity & Dimensional & Divide by & Nondimensional \\
\midrule
Position       & $\tilde{\vecr}$ [km]      & $L$     & $\vecr = \tilde{\vecr}/L$ \\
Velocity       & $\tilde{\vecv}$ [km/s]    & $V$     & $\vecv = \tilde{\vecv}T/L$ \\
Time           & $\tilde{t}$ [s]           & $T$     & $t = \tilde{t}/T$ \\
Acceleration   & $\tilde{a}$ [km/s$^2$]    & $A$     & $a = \tilde{a}T^2/L$ \\
Mass           & $\tilde{m}$ [kg]          & $m_0$   & $m = \tilde{m}/m_0$ \\
\bottomrule
\end{tabular}
\end{center}

Note the last row: mass is scaled by the \emph{initial wet mass}, not by $M$, so
$m$ is a mass fraction beginning at $1$. This is why \eqref{eq:mdot} carries no
separate $m_0$.

\paragraph{Useful reference values.} One nondimensional time unit is
$T = 4.4327$~days, so a full revolution of the primaries, $2\pi$ in these
units, is $27.851$~days. The lunar radius is
$R_M/L = 4.458264\times10^{-3}$, and the $441$~km closest approach reported in
Section~\ref{sec:illposed} is $\rho = 5.58989\times10^{-3}$---which is worth
internalizing, because a lunar \emph{impact} and a healthy $10\,000$~km flyby
differ only in the third decimal place of a state variable whose other
components are of order unity. That numerical near-degeneracy is exactly why
the constraint of Section~\ref{sec:pathcon} is needed.

\paragraph{Propulsion, which is the fiddly one.} The propulsion parameters need
care because they mix mass, force and time. Standard gravity becomes a
nondimensional number,
\[
  g_0^{\text{ND}} = \frac{9.80665~\text{m/s}^2 \cdot T^2}{1000\,L} = 3690.98 ,
\]
and then, for a thruster of maximum thrust $F$ [N] on a wet mass $m_0$ [kg] with
specific impulse $I_{sp}$ [s],
\[
  T_{\max} = \frac{F}{m_0}\cdot\frac{T^2}{1000\,L},
  \qquad
  c = \frac{I_{sp}}{T} \cdot g_0^{\text{ND}} .
\]
The first is an \emph{acceleration} at unit mass fraction, not a force; the
second is the exhaust velocity. For the $70$~mN, $900$~s, $150$~kg case used
throughout this document,
\[
  T_{\max} = 0.1756418 \;(\,= 0.4667~\text{mm/s}^2\,),
  \qquad
  c = 8.673746 \;(\,= 8.826~\text{km/s}\,),
\]
and the sanity check on the second is immediate: $I_{sp}g_0 = 900 \times 9.80665
= 8826$~m/s, recovered exactly by multiplying $c$ by $V$. Any conversion of this
kind should be checked by round-tripping one such known quantity; it is the
cheapest bug filter available.

\paragraph{Reading results back out.} The two numbers most often wanted
dimensionally are the transfer duration, $t_f T / 86400$ in days, and the
$\Delta V$, which follows from the ideal rocket equation applied to the mass
fraction: $\Delta V = c\,V \ln(1/m(t_f))$ in km/s.

\section{The two orbits}

Both endpoints are periodic solutions of the Earth--Moon circular restricted
three-body problem (CR3BP), drawn from catalogs in \code{pumpkynPie} and closed
by natural-parameter continuation before use.

\subsection{The distant retrograde orbit (departure)}

A \emph{distant retrograde orbit} (DRO) is a planar, closed orbit about the
smaller primary---here the Moon---traversed \emph{retrograde} in the rotating
frame, meaning the spacecraft circulates opposite to the Moon's motion about the
Earth--Moon barycentre. ``Distant'' refers to its size relative to the secondary:
these orbits stand well off the Moon and are strongly influenced by both bodies.

Their defining practical property is \textbf{stability}. Unlike halo and
Lyapunov orbits, which are unstable and require continual station-keeping, DROs
of moderate size are linearly stable, so a spacecraft parked in one stays there.
That makes a DRO a natural \emph{staging point}: a place to accumulate
spacecraft cheaply before dispersing them to operational orbits.

The departure orbit used here has nondimensional period $\tau_0 = 1$ and initial
state
\[
  \vecr_0 = (0.9145255582,\; 0,\; 0), \qquad
  \vecv_0 = (0,\; 0.4921644867,\; 0),
\]
in the rotating barycentric frame, nondimensionalized so that the Earth--Moon
separation, the system's mean motion, and the total mass are unity.

\subsection{The tulip orbit (arrival)}

A \emph{tulip orbit} is a multi-petal periodic orbit of the CR3BP whose ground
track traces a flower-like pattern, with $N_p$ petals arranged about the Moon
and apoapses concentrated over one lunar hemisphere. The family is indexed by
petal count $N_p$ and by a hemisphere sign; we take $N_p = 7$ and the southern
branch.

Their interest is navigational. A constellation in southern-hemisphere tulip
orbits provides sustained high-elevation geometry over the lunar south pole,
where surface operations are planned and where conventional relay architectures
give the weakest coverage. That geometric advantage is the subject of separate
work (arXiv:2607.10482); here the tulip is simply the destination.

The arrival orbit has period $\tau_f = 5 \cdot 2\pi/6$. The arrival
\emph{point} on that orbit is a design choice, not a given: we adopt the
convention of the reference implementation and select the phase whose velocity
direction is most opposed to the departure velocity, giving
\[
  \vecr_f = (0.9772215792,\; -0.0017264861,\; 0.0129258427),
\]
\[
  \vecv_f = (-0.4221963472,\; -0.9757277486,\; -0.2816581452).
\]

\section{The dynamics}

In the rotating barycentric frame, nondimensionalized, with mass ratio
$\mu = 0.012150585609624$, the Earth sits at $(-\mu,0,0)$ and the Moon at
$(1-\mu,0,0)$. Write the distances to each primary as
\[
  d = \lVert \vecr - (-\mu,0,0) \rVert, \qquad
  \rho = \lVert \vecr - (1-\mu,0,0) \rVert .
\]
The controlled equations of motion for position $\vecr$, velocity $\vecv$ and
mass fraction $m$ are
\begin{align}
  \dot{\vecr} &= \vecv, \label{eq:rdot}\\
  \dot{\vecv} &= \mathbf{g}(\vecr) + \mathbf{h}(\vecv)
                 + \frac{u\,T_{\max}}{m}\,\bal, \label{eq:vdot}\\
  \dot{m}     &= -\frac{T_{\max}}{c}\,u, \label{eq:mdot}
\end{align}
where the gravity-plus-centrifugal and Coriolis terms are
\[
  \mathbf{g}(\vecr) =
  \begin{pmatrix}
    x - (1-\mu)\dfrac{x+\mu}{d^3} - \mu\dfrac{x-1+\mu}{\rho^3}\\[2mm]
    y - (1-\mu)\dfrac{y}{d^3} - \mu\dfrac{y}{\rho^3}\\[2mm]
    \phantom{y} - (1-\mu)\dfrac{z}{d^3} - \mu\dfrac{z}{\rho^3}
  \end{pmatrix},
  \qquad
  \mathbf{h}(\vecv) = \begin{pmatrix} 2\dot y \\ -2\dot x \\ 0 \end{pmatrix}.
\]
Here $T_{\max}$ is the maximum thrust acceleration at unit mass fraction and $c$
the exhaust velocity, both nondimensional. The controls are the thrust
\emph{direction} $\bal$, a unit vector, and the \emph{throttle} $u \in [0,1]$.

Note that $m$ is a mass \emph{fraction}, starting at unity, so the thrust
acceleration $u T_{\max}/m$ grows as propellant is spent.

\section{The optimal control problem}

\begin{equation}
\begin{aligned}
  \min_{\bal(\cdot),\,u(\cdot),\,t_f} \quad & t_f \\
  \text{subject to}\quad
   & \eqref{eq:rdot}\text{--}\eqref{eq:mdot} \quad \text{on } [0,t_f], \\
   & \lVert \bal(t) \rVert = 1, \qquad u(t) \in [0,1], \\
   & \vecr(0) = \vecr_0,\quad \vecv(0) = \vecv_0,\quad m(0) = 1, \\
   & \vecr(t_f) = \vecr_f,\quad \vecv(t_f) = \vecv_f .
\end{aligned}
\label{eq:ocp}
\end{equation}

The final mass is free---minimum time does not care what it costs---and the
final time $t_f$ is itself a decision variable.

\label{sec:ocp_pmp}
\paragraph{What Pontryagin says about the answer.} Although we solve
\eqref{eq:ocp} directly, the indirect view predicts the structure and is worth
stating because we use it as a check. With costates
$(\boldsymbol{\lambda}_r, \boldsymbol{\lambda}_v, \lambda_m)$, the optimal
thrust direction is anti-parallel to the \emph{primer vector},
$\bal^\star = -\boldsymbol{\lambda}_v / \lVert \boldsymbol{\lambda}_v \rVert$,
and the throttle is bang-bang on the sign of a switching function,
\[
  S = -\frac{\lVert \boldsymbol{\lambda}_v \rVert\, c}{m} - \lambda_m,
  \qquad u^\star = \begin{cases} 1, & S < 0,\\ 0, & S > 0.\end{cases}
\]
For this problem one expects $u^\star \equiv 1$, and in every converged solution
here it is: $\min_t u(t) = 1.000000$ to six figures at every mesh density tried.
That is an \emph{empirical} property of these extremals, not a consequence of
the objective---minimum-time problems can admit genuine coast arcs when the
endpoint geometry or an active path constraint makes one optimal. We therefore
leave $u$ free so that saturation is observed rather than imposed, and the gate
of Section~\ref{sec:gate} reports it as \emph{advisory} rather than as a
required condition.

Two further conditions belong to this formulation and are stated here because
both are used later. The final mass is free, giving the transversality condition
$\lambda_m(t_f) = 0$; and the final time is free with autonomous dynamics,
giving $H(t_f) = 0$ and hence $H \equiv 0$ along the extremal.

\section{Discretization}
\label{sec:disc}

\subsection{Normalized time, so that \texorpdfstring{$t_f$}{tf} can be free}

Collocation needs a fixed integration interval, but $t_f$ is unknown. Introduce
$s = t/t_f \in [0,1]$; then for any state $X$,
\[
  \frac{dX}{ds} = t_f \, f(X,U),
\]
and $t_f$ becomes an ordinary decision variable multiplying the dynamics.

\subsection{Trapezoidal collocation}

Partition $[0,1]$ into $N$ intervals with nodes
$0 = s_1 < \cdots < s_{N+1} = 1$, $\Delta s_k = s_{k+1}-s_k$, and carry the
state $X_k = (\vecr_k, \vecv_k, m_k) \in \mathbb{R}^7$ and control
$U_k = (\bal_k, u_k) \in \mathbb{R}^4$ at every node. The dynamics are imposed
as \emph{defect} constraints,
\begin{equation}
  X_{k+1} - X_k
  - t_f\,\frac{\Delta s_k}{2}\Bigl[ f(X_k,U_k) + f(X_{k+1},U_{k+1}) \Bigr] = 0,
  \qquad k = 1,\dots,N,
\label{eq:defect}
\end{equation}
together with $\lVert \bal_k \rVert^2 = 1$, $0 \le u_k \le 1$, the boundary
conditions of \eqref{eq:ocp}, and $m_k \ge 0.05$ to keep the mass physical (a
non-strict bound: an NLP cannot impose a strict inequality). The
objective is simply $t_f$.

\subsection{Why \texorpdfstring{$t_f$}{tf} must be lifted}
\label{sec:lift}

Written as above, the single scalar $t_f$ appears in \emph{every} defect
\eqref{eq:defect}. In the Karush--Kuhn--Tucker system this produces one
\textbf{dense column}: an arrowhead structure whose factorization fills in
catastrophically.

This is not a theoretical concern. The first implementation used a scalar $t_f$
and MATLAB terminated with a fatal error immediately after IPOPT printed the
Jacobian structure, at only $N = 400$. The same hazard is documented in two
sibling campaigns: the GTO$\to$tulip solver avoids it by fixing its regularized
horizon outright, and the Earth CR3BP campaign cured it with a technique its
authors call \code{liftDL}.

The cure is to \emph{lift} $t_f$: replicate it as a per-node variable $T_k$ and
impose local continuity,
\begin{equation}
  T_{k+1} - T_k = 0, \qquad k = 1,\dots,N,
\label{eq:lift}
\end{equation}
using $T_k$ in defect $k$. The continuity constraints force all $T_k$ equal, so
the solution is unchanged, but each defect now touches only its own two nodes
and the arrowhead becomes \textbf{banded}. With the lift in place the same
problem solves in $6.8$ seconds. As a check, the spread
$\max_k |T_k - T_1|$ is reported and comes back identically zero.

\section{Fourth order: Hermite--Simpson}
\label{sec:hs}

The trapezoidal scheme of Section~\ref{sec:disc} is second-order accurate, and
Section~\ref{sec:contres} shows that is not enough on this problem at any node
count tried. The remedy applied is a fourth-order scheme at the same node count.

\paragraph{Global order and local order are different numbers, and the
diagnostic measures the second.} Hermite--Simpson has fourth-order \emph{global}
state accuracy for smooth arcs, and correspondingly a \emph{local} truncation
error of $O(h^5)$. The residual of Section~\ref{sec:contres} restarts the
integration at every node, so it measures the local error. Halving $h$ should
therefore reduce it by about $2^5 = 32$; the measured factor is $37.6$, which is
the consistency check on the implementation rather than an anomaly. Quoting the
scheme as ``fourth order'' while measuring $h^5$ is standard but worth stating
explicitly, since the two numbers otherwise look contradictory.

\subsection{The scheme}

Over interval $k$, with step $h_k = t_f \Delta s_k$, introduce an interior
midpoint state $X_{m,k}$ and midpoint control $U_{m,k}$, and impose two
conditions instead of one:
\begin{align}
  \text{interpolation:}\quad & X_{m,k} - \tfrac{1}{2}(X_k + X_{k+1})
      - \tfrac{h_k}{8}\bigl[f(X_k,U_k) - f(X_{k+1},U_{k+1})\bigr] = 0,
      \label{eq:hsinterp}\\
  \text{defect:}\quad & X_{k+1} - X_k
      - \tfrac{h_k}{6}\bigl[f(X_k,U_k) + 4 f(X_{m,k},U_{m,k})
      + f(X_{k+1},U_{k+1})\bigr] = 0. \label{eq:hsdefect}
\end{align}
Equation~\eqref{eq:hsinterp} is the Hermite cubic through the two node states
and their derivatives, evaluated at the midpoint; \eqref{eq:hsdefect} is
Simpson's rule through the three points.

\paragraph{The midpoint control is a variable, not an average.} It would be
simpler to set $U_{m,k} = \tfrac{1}{2}(U_k + U_{k+1})$. Two reasons not to.
Tying the control to a piecewise-linear family costs the scheme its fourth
order; and here the direction $\bal$ lives on the unit sphere, where the average
of two unit vectors is not a unit vector. $U_{m,k}$ therefore carries its own
norm constraint and throttle bounds.

\paragraph{Separated, not compressed---and this was learned the hard way.}
Equation~\eqref{eq:hsinterp} can be eliminated by substituting the interpolant
directly into $f$ in \eqref{eq:hsdefect}, halving the constraint count. That
form is algebraically identical, and it was implemented first. It solves at
$N = 400$ and \textbf{MATLAB terminates silently at $N = 800$}: substituting
$X_m(X, t_f)$ inside $f$ means the Lagrangian Hessian requires the chain rule
through the interpolant, the expression graph deepens with every node, and MUMPS
does not survive it.

This is the third appearance of the same failure in this repository, and the
same cure works every time: keep each constraint \emph{shallow} and let extra
variables carry the coupling. It is the reason $t_f$ is lifted
(Section~\ref{sec:lift}), the reason the earth CR3BP campaign has a routine
called \code{liftDL}, and the reason $X_m$ is a decision variable here rather
than an expression.

\paragraph{The floor gets a free upgrade.} Because $X_{m,k}$ exists as a
variable, the altitude constraint of Section~\ref{sec:pathcon} is imposed at the
midpoints as well as the nodes. That halves the unguarded span between
enforcement points. It does not eliminate it.

\subsection{What it bought}

Seeded from the indirect reference, unconstrained, at matched node counts:

\begin{table}[h]
\centering\small
\begin{tabular}{llrrr}
\toprule
scheme & form & $N$ & $t_f$ & worst position error \\
\midrule
trapezoid        & ---         & $400$  & $4.1438682$ & $1441$ km \\
Hermite--Simpson & compressed  & $400$  & $4.0173381$ & --- \\
Hermite--Simpson & separated   & $400$  & $4.6808938$ & $7.38$ km \\
\midrule
trapezoid        & ---         & $1600$ & $7.253098$  & --- \\
Hermite--Simpson & separated   & $1600$ & $\mathbf{4.0152501}$ & $\mathbf{0.0033}$ \textbf{km} \\
\bottomrule
\end{tabular}
\end{table}

\textbf{The two $N = 400$ Hermite--Simpson rows are not a typographical
error}, and an earlier version of this document listed only the first without
saying which form produced it. The compressed and separated forms are
algebraically the same problem, but the separated form carries $7N$ additional
variables and $7N$ additional equalities, so the interior-point path runs through
a different space and can converge to a different local minimum---which here it
does. Since the compressed form cannot be used beyond $N = 400$ at all, the
separated row is the one that matters; it is also the one that goes on to
certify at $N = 1600$ (Section~\ref{sec:gate}).

\subsection{A warm start must be feasible for the constraints that define it}

Lifting $X_m$ to a variable creates an obligation: it must be initialized to
satisfy \eqref{eq:hsinterp}. Seeding it with the plain midpoint average
$\tfrac{1}{2}(X_k + X_{k+1})$ leaves the starting point infeasible, and IPOPT
begins in restoration rather than in optimization.

The honest record is that this mattered less predictably than expected. Fixing
the seed left the $N = 400$ answer bit-identical at $t_f = 4.6808938$---so the
hypothesis that it explained a basin discrepancy was \emph{wrong}---while
changing the $N = 800$ answer from $4.6850952$ to $4.1628670$. The change is
correct on principle; the reasoning first offered for it was not.

\section{Adding a path constraint}
\label{sec:pathcon}

Section~\ref{sec:illposed} shows that this problem needs a minimum-altitude
constraint to be well posed. This section is about how one is added, why the
direct method makes it nearly free, and---honestly---what adding it did and did
not fix.

\subsection{The constraint, and its transcription}

The requirement is a \emph{state path constraint}: at every instant of the
transfer, the spacecraft must stay above a chosen altitude $h_{\min}$ over the
lunar surface,
\begin{equation}
  \rho(t) \;\ge\; \rho_{\min} \;=\; \frac{R_M + h_{\min}}{L},
  \qquad \forall\, t \in [0, t_f].
\label{eq:pathcon}
\end{equation}
In the transcription this becomes one algebraic inequality per enforcement
point---the $N+1$ nodes under trapezoidal collocation, and the nodes plus the $N$
interior midpoints under Hermite--Simpson, where the midpoint state exists as a
variable. We impose it in \emph{squared} form,
\begin{equation}
  \lVert \vecr_k - (1-\mu,0,0) \rVert^2 \;\ge\; \rho_{\min}^2,
  \qquad k = 1,\dots,N+1,
\label{eq:pathcon_disc}
\end{equation}
which avoids the square root---whose gradient is undefined at $\rho = 0$ and
poorly scaled near it---and keeps the constraint a smooth quadratic in the
position variables. Its gradient, $2(\vecr_k - \vecr_{\!M})$, is exact and
well conditioned everywhere the constraint could plausibly be active.

That is the entire implementation: four lines of CasADi, one extra block of
$N+1$ rows in the constraint Jacobian, no change to the objective, the dynamics,
the boundary conditions, or the solver settings.

\subsection{Is this a genuine advantage of the direct method?}

Yes---and the asymmetry is larger than it first appears.

For the \textbf{direct} method a path constraint is structurally the same kind
of object as everything else in the NLP. The problem is already a constrained
optimization; \eqref{eq:pathcon_disc} adds one sparse block of inequality rows
to a Jacobian that already carries tens of thousands. The interior-point solver handles the inequality with the same
barrier machinery it uses for $0 \le u_k \le 1$. Nothing about the formulation
changes. The switch from an unconstrained to a constrained problem cost one
optional argument in the solver's interface.

For the \textbf{indirect} method the same constraint changes the mathematical
object being solved. Under Pontryagin's principle with a pure state
inequality, one must:
\begin{enumerate}
  \item guess, \emph{a priori}, the \textbf{structure} of the solution---how
        many constrained arcs there are and in what order they appear, since
        this determines the number of unknowns in the shooting problem;
  \item derive the order $q$ of the constraint, meaning how many times
        \eqref{eq:pathcon} must be differentiated along the dynamics before the
        control appears (for an altitude constraint under thrust, $q = 2$);
  \item adjoin the constraint to the Hamiltonian with a multiplier
        $\eta(t) \ge 0$ supported only on the boundary arc, together with the
        tangency conditions $C = C^{(1)} = 0$ at entry;
  \item impose \textbf{junction conditions} at entry to the boundary arc, where
        the costate is permitted to \emph{jump}:
        $\boldsymbol{\lambda}(\tau^-) = \boldsymbol{\lambda}(\tau^+) +
        \sum_{j=0}^{q-1}\nu_j \,\partial_x C^{(j)}$, introducing further unknown
        multipliers $\nu_j$, where $\partial_x$ is a gradient with respect to the
        \emph{full} state $(\vecr,\vecv,m)$;
  \item solve the resulting multipoint boundary-value problem, whose unknown
        vector now contains the entry and exit times as well.
\end{enumerate}
The shooting problem grows, it acquires discrete structure that must be guessed
before it can be solved, and the costate trajectory is no longer continuous.
This is the same category of difficulty that makes min-fuel harder than
min-time for the indirect method---the burden of guessing a discrete structure
in advance---and here it arrives even in a min-time problem.

\paragraph{Three qualifications on the paragraph above, all of them from
review.} First, the formulation just described is the \emph{direct-adjoining}
one (Hartl, Sethi and Vickson, 1995). It is not the only convention: an
\emph{indirect-adjoining} formulation appends the $q$-th derivative $C^{(q)}$ to
the Hamiltonian and yields \emph{continuous} costates with discontinuous
derivatives instead. So the costate discontinuity above is a property of the
convention, not of the problem, and the convention must be named---as it now is.
Second, the sign restrictions are not uniform across the $\nu_j$: the multiplier
attached to the constraint itself is sign-restricted, while the tangency
multipliers $\nu_1,\dots,\nu_{q-1}$ are in general \emph{unrestricted}. Third,
the comparison should be narrowed to what it actually supports: direct
transcription makes a state-constrained problem substantially easier to
\emph{pose and numerically explore}. It is not free---the direct method still
owes continuous-time constraint verification and mesh refinement---and modern
indirect practice need not always guess arcs by hand, since continuation,
multiple shooting and boundary-arc detection all exist.

\paragraph{The caveat, and it is a real one.} The direct method's advantage is
in \emph{formulation}, not in \emph{rigour}. Constraint
\eqref{eq:pathcon_disc} is imposed only at a \textbf{finite set of points} --
the nodes, plus the interior midpoints under Hermite--Simpson. Between two
adjacent nodes the interpolated trajectory is free to dip below $\rho_{\min}$,
and nothing in the NLP forbids it. This matters here more than it usually would,
because the constraint is active precisely at periselene---where the trajectory
has its greatest curvature and where, as Section~\ref{sec:contres} shows, the
collocation is at its least accurate. A direct solution that reports
$\min_k \rho_k = \rho_{\min}$ should be read as ``the nodes clear the floor,''
not ``the trajectory clears the floor.'' The honest check is to re-integrate the
converged solution at high accuracy and measure the true minimum---which is what
the diagnostic of Section~\ref{sec:contres} does.

Against that, the indirect formulation, once its structure is guessed correctly,
enforces \eqref{eq:pathcon} \emph{exactly} on the constrained arc, because the
arc is integrated as a boundary of the feasible set rather than sampled. So the
trade is the familiar one in this repository: the direct method makes the
problem easy to \emph{pose} and gives an approximate answer; the indirect method
makes it hard to pose and gives an exact one.

\subsection{What happened when we imposed it}

The floor is exposed as \code{opts.minAltKm} in \code{run\_dro\_tulip.m}, empty
by default so that the ill-posed baseline still reproduces. Before running the
first test we recorded a prediction: with a floor in place, $t_f$ should stop
depending on $N$, and should \emph{rise} as the floor rises, a shallower flyby
being a weaker gravity assist.

\textbf{The prediction failed}, and the way it failed is informative. At
$h_{\min} = 100$~km:

\begin{center}
\begin{tabular}{rrrrl}
\toprule
$N$ & $t_f$ & achieved min.\ alt.\ [km] & max defect & status \\
\midrule
$400$  & $4.357303$  & $1048$ & $1.4\times10^{-14}$ & Solve\_Succeeded \\
$800$  & $4.724866$  & $4680$ & $3.1\times10^{-16}$ & Solve\_Succeeded \\
$1600$ & $24.65$     & $646$  & $1.2$               & Max\_Iterations \\
\bottomrule
\end{tabular}
\end{center}

Two things are visible. First, the constraint was \textbf{never active}: every
converged solution clears $100$~km by a factor of ten or more, so
\eqref{eq:pathcon_disc} holds strictly and its multipliers are zero. Second, and
despite that, the answers are \emph{different from the unconstrained ones}
($4.144$, $3.881$, $7.253$ at the same three meshes). An inactive inequality
cannot change a local minimum---but it can, and evidently did, change which
local minimum the interior-point method \emph{walks to}, because IPOPT's barrier
terms depend on the distance to every inequality including the inactive ones.

The mesh dependence therefore did not collapse. That is the correct diagnosis of
what the altitude floor does and does not buy: it removes the \emph{unbounded}
family of ever-deeper flybys, which is a genuine repair of the ill-posedness,
but it does not by itself remove the multiplicity of local minima that the
solver can land in. Those are two different pathologies, and this experiment
separated them.

\section{Solution by IPOPT}

The transcription is built with CasADi's \code{Opti} layer, which supplies exact
first and second derivatives by automatic differentiation---no finite
differencing anywhere---and hands the problem to IPOPT, an interior-point solver
for large sparse nonlinear programs, with MUMPS as the sparse linear solver.

For the \emph{trapezoidal} scheme at $N = 800$ the program carries
$7 \times 801$ state variables, $4 \times 801$ controls and $801$ lifted time
copies, against roughly $5600$ defect equations plus the norm, boundary and
continuity constraints. The separated Hermite--Simpson scheme of
Section~\ref{sec:hs} adds $7N$ midpoint states, $4N$ midpoint controls, $7N$
interpolation equalities and $N$ further norm constraints on top of that.
Tolerance is set to $10^{-10}$.

\paragraph{Verifying the answer.} A converged NLP is not by itself a solved
optimal control problem, so each solve reports: the maximum defect recomputed
from the returned trajectory; the worst deviation of $\lVert\bal_k\rVert$ from
unity; the terminal boundary error; the lifted-time spread; and
$\min_k u_k$. Typical values are a defect of $10^{-14}$ to $10^{-16}$, a unit
error of $10^{-12}$, and a terminal error of exactly zero.

\section{Lessons learned}

\subsection{We assumed no Sundman regularization was needed. That was wrong.}

Our Earth-departure campaigns cannot be solved without a Sundman
regularization---a change of independent variable $dt/d\tau = \kappa(r)$ that
stretches the mesh through perigee, where an eccentric spiral's dynamics are
violently stiff. It would have been natural to carry that machinery over.

Characterizing the reference solution first showed it is unnecessary:

\begin{table}[h]
\centering\small
\begin{tabular}{lr}
\toprule
revolutions about the Moon & $1.18$ \\
Earth distance & $0.8535$ -- $1.1122$ (never approaches Earth) \\
closest lunar approach & $0.0164$ ND $\approx 4674$ km altitude \\
\bottomrule
\end{tabular}
\end{table}

Slightly over one revolution, with no Earth passage at all. The stiffness that
forces regularization on a forty-revolution Earth spiral is simply absent, and
plain time-domain collocation is the appropriate tool. Importing the
regularization would have been machinery adopted out of habit rather than need.

The general lesson is worth stating separately: \textbf{characterize the
trajectory before choosing the transcription.} The reference solution was
available and took minutes to interrogate; that interrogation determined the
formulation.

\paragraph{This lesson was wrong, and the measurement that refutes it.} The
argument above reasons from the \emph{shape} of the reference trajectory---few
revolutions, no Earth passage---to the adequacy of a uniform mesh. That
inference does not hold, and the continuous residual of
Section~\ref{sec:contres} measures by how much. Applied to all three converged
solutions:

\begin{table}[h]
\centering\small
\begin{tabular}{rrrrr}
\toprule
$N$ & closest approach & NLP defect & worst POSITION error & worst VELOCITY error \\
\midrule
$800$  & $441$ km  & $1.4\times10^{-14}$ & $1441$ km & $1568$ m/s \\
\bottomrule
\end{tabular}
\end{table}

That row is the pathological deep flyby, and a position error of $1441$~km
against a defect of $10^{-14}$ is the whole point. \textbf{The trapezoidal
solutions at other mesh densities are not tabulated here}, because they converged
to different local minima with different final times and so cannot be read as a
refinement sequence---comparing them would be comparing three different
trajectories, not one trajectory at three resolutions. What can be said is that
no trapezoidal solve on this problem passed the accuracy gate of
Section~\ref{sec:gate} at any density tried, while the fourth-order scheme does.

So the honest statement is not the hedged one. It is: \emph{a uniform-in-time
trapezoidal mesh is inadequate for this transfer at every density tried,
including at a sane periselene.} Something is needed---a Sundman-type change of
independent variable, a mesh concentrated near closest approach, or a
higher-order scheme---and which of the three is best is an open question, not a
settled one.

Panel (c) of the diagnostic figure shows why a uniform mesh cannot be rescued
by density alone: on the $442$~km solution the grid is uniform in time to a
ratio of $1.000$, but the lunar angle swept per interval varies by a factor of
$\mathbf{353}$ between cruise and closest approach. The mesh has no idea the
flyby is happening. Adding nodes everywhere to fix an error that lives in a few
intervals is the expensive way to solve the problem, which is precisely the
argument for regularization.

The general lesson stands, but its application here was backwards: characterize
the trajectory before choosing the transcription---and then \emph{measure}
whether the transcription you chose is actually resolving it. Reasoning from
revolution count alone was not enough.

\paragraph{What was done about it.} Of the three remedies listed above, only
the third was tried, and it was sufficient: fourth-order Hermite--Simpson
(Section~\ref{sec:hs}) reaches a worst-interval error of $0.32$~km at
$N = 1600$, a regime the second-order scheme did not reach at any node count
tried. \textbf{Sundman regularization and periselene mesh concentration remain
untried, and are now optional rather than necessary}---either may reach the same
accuracy at lower $N$, but neither is on the critical path. That is worth
stating plainly, because the natural reading of the paragraphs above is that a
Sundman change of variable was required. It was not.

\subsection{Without a minimum-altitude constraint the problem is ill-posed}
\label{sec:illposed}

This is the consequential finding. Sweeping the mesh density gives:

\begin{table}[h]
\centering\small
\begin{tabular}{lrrrr}
\toprule
$N$ & $t_f$ & vs.\ reference & max defect & closest lunar approach \\
\midrule
$400$  & $4.143868$ & $+3.2\%$  & $3.3\times10^{-14}$ & $4819$ km \\
$\mathbf{800}$  & $\mathbf{3.881410}$ & $\mathbf{-3.3\%}$ & $1.4\times10^{-14}$ & $\mathbf{442}$ \textbf{km} \\
$1600$ & $7.253098$ & $+80.6\%$ & $3.9\times10^{-16}$ & $4674$ km \\
\midrule
reference (indirect) & $4.015243$ & --- & --- & $4674$ km \\
\bottomrule
\end{tabular}
\end{table}

Every one of these is machine-feasible, yet $t_f$ varies non-monotonically with
the mesh, and one of them is \emph{faster than the reference minimum}. A
minimum cannot be beaten---so the reference is a \textbf{local} minimum, not a
global one.

The final column explains how. The fast solution buys its speed with a far
deeper lunar gravity assist: a $442$ km altitude pass against the reference's
$4674$ km. Nothing in \eqref{eq:ocp} forbids this. As the flyby is drawn closer
the assist strengthens and the transfer time falls, so
\[
  \inf t_f \quad\text{is approached as the trajectory grazes the lunar surface,}
\]
and ``the minimum-time transfer'' is not a single trajectory but a
\textbf{family parameterized by flyby altitude}, bounded below only by the Moon
itself.

The problem would therefore be ill-posed as ordinarily written. It would need a
path constraint,
\[
  \lVert \vecr(t) - \vecr_{\text{Moon}} \rVert \;\ge\; \rho_{\min},
\]
with $\rho_{\min}$ set by mission design---a safe periselene altitude---and the
reported minimum time would be meaningful only relative to that choice.

\paragraph{The evidence above was retracted, and then replaced by something
stronger.} The argument as first written rested on the $442$~km solution, and
the check of Section~\ref{sec:contres} invalidates that solution: it is not a
feasible trajectory of the continuous problem, so it is not evidence that a
faster one exists. The three rows of the table are three discretizations
disagreeing, not three optima. At that point the ill-posedness was a physical
conjecture---a deeper flyby really does give a stronger assist---but not a
result.

We then believed it had become a result on the fourth-order ladder of
Section~\ref{sec:hs}. At $N = 800$, seeded from the indirect reference, the
solver returns \code{Solve\_Succeeded} with a Hermite--Simpson defect at machine
precision and \textbf{a node $719.6$~km beneath the lunar surface}. It seemed
evidence of a different kind: nothing in \eqref{eq:ocp} forbids the trajectory
from passing through the Moon, and given the chance to shorten the transfer by
doing so, the optimizer took it.

\textbf{That inference was also wrong, and for the same reason as the first
one.} The $N = 800$ solution has a worst-interval position error of
$3\,127$~km. It is not a trajectory of the continuous problem either, so a node
inside the Moon in an \emph{infeasible discrete} solution proves nothing about
the continuous one. Making the same mistake twice, one section apart, is worth
recording as its own lesson: \emph{a converged NLP is evidence about the
transcription until its continuous residual has been measured, and about nothing
else.}

What the $N = 800$ result does show is that the discrete transcription will
exploit a missing physical exclusion constraint when one is absent. That is a
statement about our formulation, and it is reason enough to impose the floor of
Section~\ref{sec:pathcon} on engineering grounds. Whether the \emph{continuous}
unconstrained problem attains its infimum remains \textbf{open}. Settling it
needs a sequence of solutions that each pass the accuracy gate, at successively
lower altitude floors, with $t_f$ falling monotonically as the floor drops. That
experiment is specified and has not been run.

\paragraph{Why the indirect method did not produce these solutions.} Two
reasons, and only the first was appreciated at the time.

First, shooting explores only what its initial costate guess can reach, so it
returns the extremal near its seed and gives no hint whether others exist.

Second: \textbf{the indirect method has no fixed mesh whose error the optimizer
can spend.} \code{tfMinProp} integrates the state and costate equations with an
adaptive-step integrator that tightens automatically where the dynamics
stiffen---which is exactly at periselene---so the particular inter-node defect
mechanism described above cannot arise. That is a real structural difference,
and it is the reverse of the path-constraint advantage of
Section~\ref{sec:pathcon}: each method is protected against the failure mode
that afflicts the other.

It should not be overstated. Adaptive integration does not make shooting immune
to near-singular dynamics, to convergence on the wrong extremal, or to failure
near a close approach; it removes one specific failure mode, not a class of
them.

\subsection{The outstanding check, run: machine-tight defects do not mean an
accurate trajectory}
\label{sec:contres}

The previous version of this document listed an open question: whether the
$442$~km solution is \emph{genuinely} feasible, given that a direct method
enforces the dynamics only at the nodes and never looks between them. That
check has now been run, and the answer is a clean no.

\paragraph{The measurement.} Take the converged solution. For each interval
$[t_k, t_{k+1}]$, start at the node state $X_k$, apply the control linearly
interpolated between $U_k$ and $U_{k+1}$, and integrate the \emph{true}
equations of motion with \code{ode113} at tolerances $10^{-11}/10^{-13}$. Then
compare the result against the node the collocation claims to reach:
\[
  R_k \;=\; \bigl\lVert\, \Phi(t_{k+1}; t_k, X_k, U(\cdot))
                          \;-\; X_{k+1} \,\bigr\rVert .
\]
$R_k$ is the local error the scheme is committing, and it is a different
quantity from the NLP defect, which measures only how well the returned numbers
satisfy the collocation \emph{equations}.

\paragraph{Two things this measurement is not.} First, the reconstructed control
$U(\cdot)$ is a \emph{convention}, not a consequence: the transcription
constrains the dynamics only at the collocation samples and says nothing about
the control between them. We use the reconstruction consistent with each
scheme's derivation---linear for trapezoid, the Lagrange quadratic through
$(U_k, U_{m,k}, U_{k+1})$ for Hermite--Simpson---and report which was used
alongside the number. Second, $R_k$ must be resolved into components before it
is given units. The state is $(\vecr, \vecv, m)$, mixing a length, a
length-per-time and a dimensionless fraction; the norm of the seven-component
difference cannot be converted to kilometres by multiplying by $L$. An earlier
version of this document did exactly that, and the resulting figures were wrong
in both directions. Position, velocity and mass errors are now reported and
gated separately.

\paragraph{The result.} On the $N = 800$, $442$~km solution:

\begin{table}[h]
\centering\small
\begin{tabular}{lr}
\toprule
NLP defect (trapezoid equations) & $1.4\times10^{-14}$ \\
worst-interval \textbf{position} error & $\mathbf{1441}$ \textbf{km} \\
worst-interval \textbf{velocity} error & $1568$ m/s \\
lunar altitude at the worst interval & $504$ km \\
\bottomrule
\end{tabular}
\end{table}

The solution satisfies its own discrete equations to fifteen digits while
departing from the actual dynamics by over a thousand kilometres on the interval
at closest approach. The $442$~km trajectory is \textbf{not physical}. It is an
artifact of a mesh too coarse to resolve the flyby it is exploiting.

\paragraph{Why this is more than one bad case.} The error is not scattered; it
is an extremely clean function of how close to the Moon the interval sits.
Fitting the residual against lunar altitude over the whole trajectory gives
\[
  R \;\propto\; \text{altitude}^{-4.2},
\]
spanning eight orders of magnitude from $10^{-8}$ at $10^5$~km down to $O(1)$ at
periselene (panel (f) of the diagnostic figure). The trapezoidal local truncation error goes as
$h^3 \lvert \dddot{x} \rvert$, and for two-body-like motion
$\lvert\dddot{x}\rvert \sim \mu^{3/2} r^{-7/2}$, which predicts an exponent near
$-3.5$.

\textbf{That is not a match, and an earlier version of this document called it
one.} A gap of $0.74$ in a log--log exponent is a factor of five across a decade
of altitude. What the comparison supports is that the residual has the steep
inverse-distance dependence the theory leads one to expect, and no more. At least
four things could account for the difference and none has been isolated: the fit
is against \emph{altitude} $r - R_M$ rather than lunar-centre radius $r$; the
step size is uniform in normalized time while $t_f$ varies between solutions; the
Coriolis, centrifugal and thrust terms are not part of the two-body heuristic;
and a single power law is being fitted across a range that is not uniformly
asymptotic.

\paragraph{What this settles, and what it does not.} It settles the question of
Section~\ref{sec:illposed}: the mesh-dependent ``faster than the reference''
solutions are not a discovery of better optima, they are quadrature error being
optimized against. A minimum-time solver handed a discretization that
under-resolves periselene will drive the trajectory into exactly the region
where its own error is largest, because that is where the objective can be
cheaply reduced. The pathology is not that the solver failed; it is that it
succeeded on the wrong problem.

Three routes to fixing it were available: a mesh concentrated near periselene, a
Sundman regularization, or a higher-order collocation scheme. Only the third was
tried, and it was sufficient---see Section~\ref{sec:hs}, where fourth-order
Hermite--Simpson reduces the worst-interval position error to $3.3$~m. The other
two remain untried and are now optional. The altitude floor of
Section~\ref{sec:pathcon} addresses a different pathology and, as shown there,
does not remove this one.

\paragraph{The transferable lesson.} \textbf{A machine-tight defect is a
statement about the discretization, not about the trajectory.}
A corollary worth naming separately, because this campaign fell for it twice:
\emph{no conclusion about the continuous problem may be drawn from a converged
NLP whose continuous residual has not been measured}---not about optimality, not
about feasibility, and not about well-posedness. The two
quantities differ here by a factor of $10^7$. Every campaign in this repository
reports defects at the $10^{-14}$ level as evidence of a good solve; that
evidence is necessary and nowhere near sufficient, and the continuous residual
is cheap enough---one \code{ode113} call per interval---that there is no excuse
for not measuring it.

\section{Certification: what separates a converged NLP from an answer}
\label{sec:gate}

Every solve in this campaign reported a defect of $10^{-14}$ and every early one
was wrong. A gate was therefore built that a solve must pass before its $t_f$ is
quoted. Two of its conditions decide the verdict; the rest are hygiene.

\paragraph{G1, local accuracy.} The worst continuous-time local error of
Section~\ref{sec:contres}, resolved into a \textbf{position} error in kilometres
and a \textbf{velocity} error in metres per second, each with its own tolerance
(default $1$~km and $1$~m/s). Expressing the tolerances dimensionally is
deliberate: $10^{-6}$ is not obviously good or bad, whereas ``$1441$~km of
position error'' is self-evidently disqualifying. Resolving the components is
equally deliberate, and was forced on us: the first version of this gate took the
norm of the full seven-component state difference and multiplied it by $L$, which
is not a position error and cannot be read as one.

\paragraph{G1b, global accuracy.} G1 restarts at every node, so it is a maximum
over \emph{local} errors and could in principle understate what accumulates over
a whole trajectory. G1b instead flies the reconstructed control \emph{once},
from the initial node to the end, and asks where the spacecraft actually
arrives. On the certified solution the maximum local position error is $3.3$~m,
the sum of the local errors is $26.8$~m, and the end-to-end propagated error is
$42$~m---so the errors do accumulate, mildly, but nothing like the exponential
amplification a $1600$-interval integration through an unstable three-body field
might have produced. That was the expected failure mode and it did not occur.

\paragraph{G2, agreement.} DRO$\to$tulip is the one problem in the catalog where
an independently converged \emph{indirect} solution already exists. A working
direct solver must reproduce it: relative error in $t_f$ below $10^{-4}$. This
is the sharper of the two tests, and it is the reason this transfer was chosen
to begin the programme. A direct method that cannot recover a known answer where
one exists should not be believed where none does.

\paragraph{What ``certified'' means here, and what it does not.} It means the
transcription is accurate to the stated tolerances and, where a reference exists,
reproduces it. It does \textbf{not} mean the solution is a global minimum, and it
does not mean the underlying optimal control problem is well posed---the
certified solution below solves the unconstrained problem, whose well-posedness
Section~\ref{sec:illposed} leaves explicitly open. ``Certified reproduction of
the reference local extremal'' would be the fully precise phrase.

\paragraph{G3--G9, hygiene.} NLP defect, control unit-norm error, terminal
boundary error, lifted-time spread, the Hermite interpolation residual, and---when
a floor is imposed---whether the \emph{propagated} trajectory honours it between
nodes and not merely at them. These are necessary and jointly insufficient, and
they are reported separately so that a solve passing all of them while failing G1
is visible for what it is. Throttle saturation is reported as \textbf{G6*},
advisory only: as noted in Section~\ref{sec:ocp_pmp}, full thrust is an empirical
property of these extremals rather than a condition minimum-time problems must
satisfy, so it must not gate a verdict.

\paragraph{What G2 does when a floor is imposed: it switches itself off.} The
reference comes from \code{tfMin}, which has no path-constraint capability, so
for a constrained problem there is no indirect reference at all and comparing
against the unconstrained one would be worse than useless. The gate detects an
imposed floor and disables G2, saying so in its report. \textbf{No constrained
solution has yet been certified}---the constrained problem is formulated, and
that is all.

\paragraph{The instrument has its own test.} Since G1 now decides whether work is
publishable, the residual engine is itself tested: on a trajectory built by
high-accuracy integration it must report accuracy, not scheme error, and it must
detect an injected error at about the injected size. It returns
$8.3\times10^{-16}$ on the former and $1.0\times10^{-5}$ on an injected
$10^{-5}$. Three of three pass. The large residuals it reports elsewhere are
therefore real.

\subsection{The certified solution}

\begin{table}[h]
\centering\small
\begin{tabular}{llrrl}
\toprule
& gate & value & tolerance & \\
\midrule
G1   & local POSITION accuracy, worst interval & $0.0033$ km & $1.0$ km & PASS \\
G1v  & local VELOCITY accuracy, worst interval & $0.00084$ m/s & $1.0$ m/s & PASS \\
G1b  & GLOBAL POSITION accuracy, flown end to end & $0.0418$ km & $1.0$ km & PASS \\
G1bv & GLOBAL VELOCITY accuracy, end to end & $0.00185$ m/s & $1.0$ m/s & PASS \\
G2   & agreement with the indirect $t_f$ & $1.900\times10^{-6}$ & $10^{-4}$ & PASS \\
G3   & NLP defect & $2.91\times10^{-16}$ & $10^{-9}$ & PASS \\
G4   & control unit-norm error & $2.22\times10^{-16}$ & $10^{-8}$ & PASS \\
G5   & terminal boundary error & $0$ & $10^{-9}$ & PASS \\
G6*  & throttle saturation (advisory) & $1.000000$ & $1$ & PASS \\
G8   & lifted-time spread & $0$ & $10^{-10}$ & PASS \\
G9   & Hermite interpolation residual & $4.44\times10^{-16}$ & $10^{-9}$ & PASS \\
\bottomrule
\end{tabular}
\end{table}

\noindent Hermite--Simpson, $N = 1600$: $t_f = 4.0152501$ against the indirect
reference's $4.0152425$, agreeing to five significant figures, with a
worst-interval position error of $3.3$~m, an end-to-end propagated position error
of $42$~m, and a closest approach of $4673$~km.
\textbf{This is the first certified direct solution of this transfer, and the
first time in this repository that a direct and an indirect method have been
shown to agree numerically rather than merely to be two formulations of the same
problem.} Two cautions on the strength of that claim. The agreement is to
\emph{five} significant figures, not six. And on its own it is agreement on
$t_f$ and on the endpoints---which the boundary conditions pin in any case. The
genuine trajectory- and costate-level comparison is the campaign's actual
objective, and it has since been carried out: see
Section~\ref{sec:costatecmp}, where the two methods are shown to agree on the
state history to $0.4$~km and on the costate field to $6\times10^{-4}$ degrees.

\subsection{The ladder is not monotone, and the middle rung is the lesson}

\begin{table}[h]
\centering\small
\begin{tabular}{rrrrl}
\toprule
$N$ & $t_f$ & rel.\ error & worst POSITION error & min.\ node altitude \\
\midrule
$400$  & $4.6808938$ & $1.66\times10^{-1}$ & $7.38$ km & $4818$ km \\
$800$  & $4.1628670$ & $3.68\times10^{-2}$ & $3127$ km & $\mathbf{-719.6}$ \textbf{km} \\
$1600$ & $4.0152501$ & $1.90\times10^{-6}$ & $\mathbf{0.0033}$ \textbf{km} & $4673$ km \\
\bottomrule
\end{tabular}
\end{table}

A negative altitude is a node beneath the lunar surface---and the same row has a
$3127$~km position error, so it is not a trajectory at all. Read together, those
two numbers are the reason Section~\ref{sec:illposed} withdraws the
ill-posedness claim rather than resting on this row: an infeasible discrete
solution placing a node inside the Moon tells us about our transcription, not
about the continuous problem.

The row does carry a lesson the certified row does not: \textbf{refinement is not
monotone.} Going from $N = 400$ to $N = 800$ made the answer worse by every
measure, and only at $N = 1600$ did the sequence land on the reference. A
practitioner reading off a single mesh density---any single density here,
including $N = 800$ with its machine-precision defect---would have been misled.
The gate exists precisely because ``it converged'' does not distinguish these
three rows.

\section{The comparison itself: the two methods find the same extremal}
\label{sec:costatecmp}

Everything to this point certifies the direct solution against the
\emph{dynamics}. This section certifies it against the \emph{other method},
which is what DRO$\to$tulip was chosen for: it is the one problem in the catalog
where an independently converged indirect solution exists, so the covector
mapping can be tested against a second opinion instead of against itself.

\subsection{Costates: mapping, and results}

For defect constraints $D_k = 0$ with multipliers $\nu_k$, stationarity of the
Lagrangian with respect to the interior states is exactly the discrete adjoint
recursion, so $\nu_k \to \lambda(t_k)$ as the mesh refines, with no $h$-scaling
for this defect form (Hager 2000). Two ambiguities remain and are
\emph{measured rather than assumed}: the sign convention is resolved against the
solution's own thrust direction---a primal quantity, free of any dual
convention---and the scale by comparing magnitudes, with the real test being
whether the fitted factor is \emph{constant} in time.

\begin{table}[h]
\centering\small
\begin{tabular}{lr}
\toprule
primer from duals vs the solution's own thrust & max $1.2\times10^{-6}$ deg \\
direct vs indirect $\boldsymbol{\lambda}_v$ direction & median $0.0006$ deg, max $0.033$ deg \\
direct vs indirect $\boldsymbol{\lambda}_r$ direction & median $0.0006$ deg, max $0.086$ deg \\
scale factor, $\boldsymbol{\lambda}_v$ & $0.999991$ (CoV $4.1\times10^{-5}$) \\
scale factor, $\boldsymbol{\lambda}_r$ & $0.999995$ \\
terminal covector (Hager) vs indirect $\lambda(t_f)$ & $4.96\times10^{-4}$ relative \\
transversality $\lambda_m(t_f)$ & $8.2\times10^{-5}$ \\
Hamiltonian, median $\lvert\lambda^{\mathsf T}f + 1\rvert$ & $9.8\times10^{-6}$ \\
\bottomrule
\end{tabular}
\end{table}

Three results are stronger than one would design for. The primer alignment is
\emph{exact} to $10^{-6}$ degrees, because KKT stationarity for the midpoint
control is literally the PMP minimum condition
$\bal = -\boldsymbol{\lambda}_v/\lVert\boldsymbol{\lambda}_v\rVert$. The scale
factor is $1.000$, not merely constant---the two methods land in the \emph{same}
normalization, which nothing required. And $\boldsymbol{\lambda}_r$ matches too,
which is the sterner test: it couples to the trajectory only indirectly and is
the classic weak spot of shooting.

The terminal covector is read from the terminal \emph{boundary} multipliers
rather than extrapolated: stationarity with respect to $X_{N+1}$ gives
$\lambda(t_f) = -\nu_\psi$ to leading order, the Hager terminal correction.

\subsection{State: the raw difference is phasing, not shape}

\begin{table}[h]
\centering\small
\begin{tabular}{lr}
\toprule
max position difference, matched \emph{absolute} time & $3.311$ km \\
max position difference, matched \emph{fractional} time & $\mathbf{0.396}$ \textbf{km} \\
median, fractional & $0.230$ km \\
max velocity difference & $0.341$ m/s \\
max mass-fraction difference & $3.2\times10^{-10}$ \\
thrust direction difference & median $5.3\times10^{-4}$ deg \\
\bottomrule
\end{tabular}
\end{table}

The two solutions differ in $t_f$ by $7.63\times10^{-6}$~ND $= 2.92$~s. At the
trajectory's speed range ($0.143$--$1.135$ km/s) that offset alone predicts an
along-track displacement of $0.42$--$3.32$~km---and the observed maximum at
matched absolute time is $3.311$~km, the top of that range. Comparing at matched
fractional time $s = t/t_f$ removes the phasing and leaves $0.396$~km of genuine
shape difference over an $89{,}000$~km, $17.8$-day transfer, at the level of the
remaining discretization error rather than of any disagreement between methods.

\subsection{A third misattribution, caught by its own test}

The Hamiltonian condition $\lambda^{\mathsf T}f \equiv -1$ holds to
$10^{-5}$ over the interior but deviates by $O(1)$ on the last few intervals. We
first attributed that to a terminal covector-mapping artifact. \textbf{The
explanation is wrong}, and the test that refutes it is one line: on this
transfer the closest lunar approach occurs at interval $1599$ of $1600$, so
``the boundary'' and ``periselene'' are the same place. Fitting
$\lvert\lambda^{\mathsf T}f+1\rvert$ against altitude on the first $1560$
intervals gives altitude$^{-4.7}$, and the last eight intervals sit at
$1.1\times$ that prediction---no excess. The deviation is ordinary truncation
error at periselene, which happens to sit at $t_f$.

That is the third time this document has attributed an effect to the wrong
cause by failing to separate two confounded variables. The comparison figures
therefore carry lunar altitude on a second axis in every panel that varies with
time, so the confound is visible rather than remembered.

\subsection{Figures}

\code{results/dvi\_N1600\_state.png} overlays the two trajectories and shows the
position difference both raw and phasing-removed, the velocity, mass and
thrust-direction differences, and the difference against altitude.
\code{results/dvi\_N1600\_costate.png} overlays $\boldsymbol{\lambda}_r$,
$\boldsymbol{\lambda}_v$ and $\lambda_m$ (direct solid, indirect dashed---they
are indistinguishable, including the large $\boldsymbol{\lambda}_r$ excursions),
with the direction agreement, the scale factor, and the Hamiltonian condition.
Generator: \code{viz/plot\_direct\_vs\_indirect.m}; the numbers come from
\code{certify/costate\_compare.m}.

\subsection{What this closes, and what it opens}

The claim ``the two methods agree,'' earlier reduced to $t_f$ plus pinned
endpoints, is now trajectory- and costate-level: same state history, same primer
field, same costate directions, same normalization. It also demonstrates the
mechanism the costate-catalog programme needs---a direct solve yields costates
accurate to $\sim10^{-3}$ degrees against a converged indirect answer, far
better than any hand-built guess---validated here on the one problem where the
validation was possible, for use on the problems where it is not.

\section{Diagnostics and visualization}

Stage 4 of \code{run\_dro\_tulip.m} is optional and off by default. It runs on
one chosen solve from the sweep and produces the following.

\paragraph{The six-panel diagnostic} (\code{opts.plots = true}), ordered by what
a reader should check first:
\begin{enumerate}
  \item[(a)] the trajectory, Moon-centred, with the lunar disc and the altitude
        floor drawn to scale---so that whether the constraint bites is visible
        rather than inferred;
  \item[(b)] lunar altitude against time, on a log axis, with the floor marked;
  \item[(c)] mesh spacing, measured two ways: $\Delta t$ per interval, and the
        lunar angle swept per interval. The first is flat by construction, the
        second is the one that matters, and their disagreement is the
        Sundman argument of the qualification above;
  \item[(d)] the continuous-time residual $R_k$ of Section~\ref{sec:contres},
        against the NLP defect drawn as a dashed line for scale;
  \item[(e)] the throttle, which for minimum time should saturate at unity
        throughout---anything else is a finding;
  \item[(f)] residual against lunar altitude, the panel that exhibits the
        power law.
\end{enumerate}

\paragraph{Two movies} (\code{opts.movies}), both written at $1280\times720$
so that H.264 does not shear the frames:
\begin{itemize}
  \item \code{'traj'}: a Moon-centred rotating-frame movie in base MATLAB only,
        with the departure DRO and arrival tulip for context, the lunar disc and
        altitude floor to scale, a fading tail behind the spacecraft, and a live
        readout of elapsed time and lunar altitude. This one always runs.
  \item \code{'scene'}: the \code{pumpkynPie}-lit scene---starfield, shaded
        Moon, Earth---driven by \emph{our} direct solution rather than the
        indirect one, built by calling \code{pumpkynPie.plot.SatelliteAnimator}.
        Nothing is written into \code{pumpkyn} or \code{pumpkynPie}. If that
        class is not on the path the function returns cleanly and the caller
        falls back to \code{'traj'}.
\end{itemize}
In both, the path between collocation nodes is drawn by spline interpolation, at
six samples per interval. That is a \emph{rendering} choice and is stated as
such: it makes the curve smooth without asserting that the sub-node path is
dynamically correct---which, per Section~\ref{sec:contres}, near periselene it
emphatically is not.

\section*{Reproducing this}
\addcontentsline{toc}{section}{Reproducing this}

\begin{verbatim}
    addpath('orbit_transfer/DRO_tulip/direct');

    R = run_dro_tulip();                        % the baseline record

    R = run_dro_tulip(struct( ...               % with everything on
            'scheme',   'hermite-simpson', ...  %   4th order; needed for G1
            'N',        1600, ...               %   where the gate passes
            'minAltKm', 500, ...                %   altitude floor, km
            'certify',  true, ...               %   G1/G2 pass-fail gate
            'plots',    true, ...               %   six-panel diagnostic
            'movies',   'both'));               %   'traj' | 'scene' | 'both'
\end{verbatim}

The first call regenerates the reference characterization, the mesh sweep and
the lunar-approach table. Requires \code{pumpkynPie} and \code{pumpkyn}---which
must be updated together---and CasADi 3.7.0. Leaving \code{minAltKm} empty is
what reproduces the ill-posed baseline of Section~\ref{sec:illposed}, so the
default is deliberately the unconstrained one---but note
Section~\ref{sec:gate}: that default problem has no minimum, and only the
fourth-order scheme at $N = 1600$ passes the gate.

\paragraph{What remains open.} The accuracy ladder should be redone with the
floor imposed, since the unconstrained problem is not a valid testbed for it.
$N = 3200$ was never run---the ladder was stopped one rung early. And G2 needs
splitting once a floor is in place, because \code{tfMin} has no path constraint
and so supplies no reference for the constrained problem: transcription fidelity
(fix $t_f$, solve for feasibility, require a match) and optimizer quality (free
$t_f$ with a floor) are different questions and deserve different gates.

\begin{tabular}{ll}
\toprule
front door & \code{DRO\_tulip/direct/run\_dro\_tulip.m} \\
pumpkyn-style companion & \code{DRO\_tulip/direct/run\_dro\_tulip\_ps.m} \\
endpoints & \code{DRO\_tulip/direct/lib/dro\_tulip\_endpoints.m} \\
solver & \code{DRO\_tulip/direct/lib/casadi\_mintime\_dro.m} \\
gate & \code{DRO\_tulip/direct/certify/certify\_dro\_mintime.m} \\
residual engine & \code{DRO\_tulip/direct/certify/dro\_residual.m} \\
its self-test & \code{DRO\_tulip/direct/certify/tests/test\_dro\_residual.m} \\
diagnostics & \code{DRO\_tulip/direct/viz/plot\_dro\_diagnostics.m} \\
movies & \code{DRO\_tulip/direct/viz/dro\_\{traj,scene\}\_movie.m} \\
running record & \code{DRO\_tulip/FINDINGS.md} \\
indirect reference & \code{pumpkynPie/demos/lowThrustDRO2Tulip.m} \\
\bottomrule
\end{tabular}

\end{document}
